%% Supplementary Material for:
%% "Closing the Security Activation Gap: A Security-by-Design Methodology for IoT
%%  Ecosystems and Its Empirical Validation in a DC Smart Microgrid"

\documentclass[12pt]{article}

\usepackage[a4paper,margin=2.5cm]{geometry}
\usepackage{amsmath,amssymb,amsfonts}
\usepackage{booktabs}
\usepackage{array}
\usepackage{xcolor}
\usepackage[utf8]{inputenc}
\usepackage[T1]{fontenc}
\usepackage[colorlinks=true,linkcolor=blue,citecolor=blue,urlcolor=blue]{hyperref}
\usepackage[authoryear,round]{natbib}

\title{Supplementary Material:\\[4pt]\large Closing the Security Activation Gap: A Security-by-Design Methodology for IoT Ecosystems and Its Empirical Validation in a DC Smart Microgrid}
\author{}
\date{}

\begin{document}
\maketitle
\renewcommand{\thetable}{S\arabic{table}}
\setcounter{table}{0}

\section*{S1. Methodology Parameterized Specification: SMG Case Study}
\label{app:stages}

This document provides the full parameterized specification of all four phases for the SMG case study, complementing the condensed traceability in the main text's Tables~10 and~11. Each row in Tables~S1--S4 corresponds to a methodology parameter defined in the methodology's process tables. Codes with an additional digit (e.g., 4.1.2.1) refine a stage-level process (as cited in the main text) into its implementation-specific components; the parent code is the methodology-level process that those refinements collectively satisfy.

\begin{table}[htbp]
    \centering
    \caption{Phase A, Stage~1: Organizational aspects --- SMG application.}
    \label{tab:app-stage1}
    \small
    \renewcommand{\arraystretch}{1.15}
    \begin{tabular}{
        >{\centering\arraybackslash}p{1.3cm}
        >{\raggedright\arraybackslash}p{3.7cm}
        >{\raggedright\arraybackslash}p{10.0cm}
    }
        \toprule
        \textbf{Process ID} & \textbf{Parameter} & \textbf{SMG Value} \\
        \midrule
        1.1   & Problem statement & Manage distributed DC renewable energy generation, storage, and consumption across 2 co-located laboratory nodes; provide global coordination and real-time monitoring without centralized national grid \\
        1.1.1   & Cost-benefit justification & Energy losses (imbalanced generation/demand) eliminated; TRL~3 budget: 2\texttimes{} ESP32 NodeMCU + RPi4 + commercial solar adapters + AGM batteries + IBFIA SSRs \\
        1.1.2   & Entities involved & 5: IoT nodes (2\texttimes{} ESP32), IoT server (RPi4), users (operators), administrators (researcher), DC bus (physical) \\
        1.3.1   & Development iteration (TRL) & ITER~1 target: TRL~3 (proof of concept with physical elements, controlled laboratory) \\
        1.4.1   & Criticality classification & Critical infrastructure --- electrical energy; high-security tier selected (Stage~1.4.1) \\
        1.4.2   & Technological sovereignty & No cloud services; all computation on ESP32 (local) and RPi4 (global); open-source software only \\
        1.4.3   & Data sovereignty & Self-hosted RPi4 server; no third-party cloud; all sensor data stored locally \\
        1.4.4   & Access control policy & Token-based; nodes: POST sensor data only; users: GET monitoring; admins: full configuration and entity management \\
        1.4.5   & Communication policy & Mandatory application-layer encryption (AES-256-GCM); TCP over WiFi~802.11 b/g/n; no plaintext data channel \\
        1.4.6   & Applicable guidelines & Secure-by-Design principles \citep{ukgov2024sbdprinciples}; NIST IoT security \citep{NIST_IoT_2020}; OWASP IoT Top~10 \\
        \bottomrule
    \end{tabular}
\end{table}

\begin{table}[htbp]
    \centering
    \caption{Phase B, Stages~2--3: Architecture and service design --- SMG application.}
    \label{tab:app-stage23}
    \small
    \renewcommand{\arraystretch}{1.15}
    \begin{tabular}{
        >{\centering\arraybackslash}p{1.3cm}
        >{\raggedright\arraybackslash}p{3.7cm}
        >{\raggedright\arraybackslash}p{10.0cm}
    }
        \toprule
        \textbf{Process ID} & \textbf{Parameter} & \textbf{SMG Value} \\
        \midrule
        2.1.4   & Communication topology & Centralized (hub-and-spoke): both nodes communicate exclusively with primary server; no node-to-node direct communication \\
        2.1.5   & Communication infrastructure & WiFi~802.11 b/g/n; ESP32 built-in; TCP sockets; RPi4 configured as WiFi AP (SSID: SMG\_Primary; NetworkManager hotspot, 10.42.0.0/24 with gateway 10.42.0.1) \\
        2.1.7   & Processing topology & Fog: local 9-rule fuzzy EMS on each ESP32; global 27-rule security-aware fuzzy EMS on RPi4; commercial solar adapter handles droop and battery-charging voltage regulation autonomously \\
        2.1.8   & Entity hierarchy & IoT Nodes $<$ Users $<$ Administrators \\
        2.2.2   & Authentication mechanism & Nonce-bound asymmetric (ECDSA) device proof + per-node symmetric (SHA-256 hash-reconstruction) server proof, both PUF-seeded, for IoT nodes; plain SHA-256 hash for users/admins (not a salted, slow key-derivation function; see main text Limitations) \\
        2.2.3   & Root of Trust process & PUF-seeded scheme (\texttt{esp32\_puflib}): enrollment via RTC-SRAM measurement (admin workstation) \textrightarrow{} non-secret mask/ECC helper data to NVS (USB serial) \textrightarrow{} HKDF-derived EC key pair \textrightarrow{} device public-key registration and per-node server-verifier provisioning \textrightarrow{} per-boot response reconstruction and mutual identification \\
        2.2.4   & Secure communications & AES-256-GCM (ESP-IDF \texttt{mbedtls}/\texttt{esp\_aes}) with 12-byte random nonces per message, keyed from the PUF-derived session key; no message-level cryptographic chain (design trade-off, see main text Limitations); no TLS certificate management needed on constrained devices \\
        3.1.1   & Services defined & 4: sensor data ingestion (nodes\textrightarrow{}server), EMS coordination (server\textrightarrow{}nodes), authentication API (all entities), monitoring dashboard (server\textrightarrow{}users) \\
        3.2.1   & Access control & RBAC via session token type; API endpoints enforce entity class; nodes POST-only to own data endpoint \\
        3.2.2   & Identity storage & ESP32: NVS flash partition (separate from firmware, survives updates, JTAG-disabled for TRL~5+); Server: PostgreSQL-equivalent for SMG + \texttt{node\_registry.json} for provisioning records \\
        \bottomrule
    \end{tabular}
\end{table}

\begin{table}[htbp]
    \centering
    \caption{Phase C, Stages~4--5: Implementation technologies and functional elements --- SMG application.}
    \label{tab:app-stage45}
    \small
    \renewcommand{\arraystretch}{1.15}
    \begin{tabular}{
        >{\centering\arraybackslash}p{1.3cm}
        >{\raggedright\arraybackslash}p{3.7cm}
        >{\raggedright\arraybackslash}p{10.0cm}
    }
        \toprule
        \textbf{Process ID} & \textbf{Parameter} & \textbf{SMG Value} \\
        \midrule
        4.1.1   & TRL~3 criteria met & Physical node operational; mutual authentication functional; all applicable functional tests passed (Table~15, main text). PWM energy control and sensor-data logging are specified by the physical hardware (Table~S3, process 5.1.2) but not yet exercised by the current C/ESP-IDF firmware, which implements only identity/RoT/secure-communication (see main text Limitations) \\
        4.1.2.1   & Node platform & ESP32 + C/ESP-IDF (v5.4.2, per the flashed firmware build log): built-in WiFi, SHA-256 HW acceleration, NVS; \texttt{esp32\_puflib} \citep{stanicek2022esp32puflib} linked directly as a native ESP-IDF component (no language-runtime wrapper needed); identity/RoT/communication firmware (\texttt{smg\_puf\_rot\_node}) builds cleanly under ESP-IDF; flashed, enrolled, and bench-validated end-to-end on SMG\_NODE\_01 ($N=30$ real hard-reset trials, 30/30 success) --- SMG\_NODE\_02 bench status still pending \\
        4.1.2.2   & Cryptography & HKDF-SHA256 (RFC~5869) derives a P-256 (secp256r1) key pair from the PUF response (\texttt{esp32\_puflib}); device authenticates via ECDSA signature over a server-issued nonce (\texttt{mbedtls}), verified against the registered public key; server authenticates via a per-node SHA-256 hash-reconstruction proof; session-key derivation $Key_{ab} = \text{SHA-256}(V_B\|N_A\|N_B)$ (\texttt{rot\_session.c}); AES-256-GCM via ESP-IDF \texttt{mbedtls}/\texttt{esp\_aes} \\
        4.1.3   & Communication protocol & TCP over WiFi~802.11 b/g/n; application-layer nonce-based mutual authentication protocol (implementation-specific Stage~4 choice, presented as part of this work's own case study) \\
        4.2.1   & Vulnerability mitigations & ESP-IDF v5.4.2 pinned (per the flashed firmware build log); WiFi credentials in NVS; all AES-GCM nonces and challenge values from \texttt{esp\_fill\_random()} (ESP-IDF HW RNG); no persistent device secret exists to protect, since the device's private key is derived fresh each boot and never stored; JTAG disable deferred to TRL~5+ \\
        5.1.1   & Primary agent (RPi4) & Python: \texttt{server.py}, \texttt{auth\_service.py}, \texttt{ems\_coordinator.py} (27-rule security-aware fuzzy), \texttt{data\_store.py}; \texttt{puf\_server.py} + \texttt{rot\_challenge\_service.py} (PUF-based mutual authentication challenge service, separate port) \\
        5.1.2.1   & Node power supply & DC-DC buck converter (12V\textrightarrow{}5V) per node from local battery; isolates ESP32 from shared DC bus; $\approx$90\% efficiency vs.\ $\approx$27\% for LDO at 12\,V (avoids $\approx$3.8\,W peak dissipation) \\
        5.1.2.2   & Identity storage & ESP32 NVS: \texttt{node\_id}, \texttt{node\_uuid}, PUF stable-bit mask, PUF ECC helper data, \texttt{wifi\_ssid}, \texttt{wifi\_pass} --- none present in \texttt{config.json} or firmware binary; no raw key or PUF response is ever written to NVS \\
        5.1.2.3   & Node firmware modules & C/ESP-IDF application \texttt{smg\_puf\_rot\_node} (\texttt{app\_main.c}, \texttt{rot\_identity.c/h}, \texttt{rot\_session.c/h}, \texttt{secure\_channel.c/h}, \texttt{net\_msg.c/h}, \texttt{wifi\_station.c/h}): identity (P1 enrollment, P2 mutual authentication challenge) and communication (P3, AES-256-GCM) modules; sensor/EMS modules not yet implemented (see main text Limitations) \\
        5.2.1   & Secure development & NVS stores only non-secret mask/ECC helper data, not a persistent key; \texttt{esp\_fill\_random()} nonces (ESP-IDF HW RNG); no credentials in source \\
        \bottomrule
    \end{tabular}
\end{table}

\begin{table}[htbp]
    \centering
    \caption{Phase D, Stages~6--7: Deployment, operation, and maintenance --- SMG application.}
    \label{tab:app-stage67}
    \small
    \renewcommand{\arraystretch}{1.15}
    \begin{tabular}{
        >{\centering\arraybackslash}p{1.3cm}
        >{\raggedright\arraybackslash}p{3.7cm}
        >{\raggedright\arraybackslash}p{10.0cm}
    }
        \toprule
        \textbf{Process ID} & \textbf{Parameter} & \textbf{SMG Value} \\
        \midrule
        6.1.3   & RoT execution & PUF-based scheme, executed and validated on SMG\_NODE\_01 (1/2 nodes; SMG\_NODE\_02 pending): admin runs \texttt{enroll\_puf()} \textrightarrow{} mask/ECC helper data written to ESP32 NVS \textrightarrow{} server verifier registration via \texttt{register\_puf\_node.py} \textrightarrow{} per-boot reconstruction and nonce-bound mutual identification confirmed over $N=30$ real hard-reset trials (30/30 success) \\
        6.2.1   & Functional test coverage & 8 tests: boot (FT-01), WiFi from NVS (FT-02), P2 handshake (FT-03), P3 session (FT-04), P4 chain (FT-05), fuzzy EMS (FT-06), ADC saturation (FT-07), degraded mode (FT-08); 100\% pass rate \\
        6.2.2   & Tamper verification & NVS parameter tampering confirmed as authentication failure (single targeted negative test; full penetration testing deferred to ITER~2) \\
        7.1.1   & Monitoring (Stage~7.1.1) & Per-node real-time: $V_{bus}$, $I_{bus}$, $P_{bus}$, $V_{gen}$, $I_{gen}$, $P_{gen}$, $D_{pwm}$; authentication status; session history (evidentiary log, no cryptographic chain, see main text Limitations); heap free \\
        7.1.4   & Maintenance cycles & One maintenance moment (SMG\_NODE\_01): an enrollment-state-confusion bug causing an unbounded re-enrollment loop on reconstruction retry, identified and fixed via code-level revision (M2) prior to the $N=30$ bench trial. M1 and M3 not triggered at TRL~3; SMG\_NODE\_02 maintenance-cycle evidence remains pending. \\
        7.2.1   & Anomaly detection & Per-node trust score (float 0.0--1.0): initialized 0.80 on successful P2 identification; $-0.30$ on P3 authentication-tag failure; $-0.40$ on $V_{bus}$ anomaly; UNTRUSTED ($<0.30$) $\Rightarrow$ TCP dropped; recovery requires re-auth \\
        7.2.2   & Update elements & OTA update requires admin token + firmware signature (deferred to ITER~2) \\
        7.2.3   & Incident response & Auth failures logged; $>3$ consecutive failures $\Rightarrow$ quarantine \\
        7.2.4   & Security-aware EMS & Global fuzzy EMS (27 rules) on RPi4 accepts the per-node trust score from 7.2.1 as a third input alongside voltage tracking error and bus current, producing security-aware energy management without a separate intrusion response layer \\
        \bottomrule
    \end{tabular}
\end{table}

\bibliographystyle{elsarticle-harv}
\bibliography{references}

\end{document}
